\documentclass[11pt]{article}
\usepackage{preamble}

\newcommand{\IconCheck}{[CHECK]}
\newcommand{\IconMic}{[MIC]}
\newcommand{\IconClipboard}{[CLIPBOARD]}
\newcommand{\lessonrow}[2]{\textbf{#1} & #2 \\ \hline}

\newcommand{\phasetag}[1]{%
  {\small\itshape Tag: #1\par}
}

\newcommand{\phaseitems}[1]{%
  \begin{itemize}[leftmargin=1.5em,itemsep=2pt,topsep=2pt]
  #1
  \end{itemize}%
}

\newcommand{\teacheranswerbox}[2]{%
  \begin{callout}{#1}{calloutgreen}
  #2
  \end{callout}
}

\newcommand{\placeholdertext}[1]{%
  {\small\itshape #1\par}
}

\newcommand{\chemistfigure}{%
\begin{center}
\begin{tikzpicture}[font=\small, line width=0.9pt]
\node[
  draw,
  rounded corners,
  fill=frameshade,
  minimum width=2.0in,
  minimum height=1.45in,
  align=center
] (chem) at (0,0) {\textbf{Chemist}\\pouring liquid\\into a container};
\node[
  draw,
  rounded corners,
  fill=calloutyellow,
  minimum width=2.0in,
  minimum height=0.62in,
  align=center
] (six) at (-2.45,-1.45) {\(x\) gallons of\\6\% alcohol};
\node[
  draw,
  rounded corners,
  fill=calloutpeach,
  minimum width=2.1in,
  minimum height=0.62in,
  align=center
] (two) at (2.45,-1.45) {50 gallons of\\2\% alcohol};
\draw[->, tealheader] (six.north east) .. controls (-1.6,-0.7) and (-0.8,-0.35) .. (chem.south west);
\draw[->, tealheader] (two.north west) .. controls (1.6,-0.7) and (0.8,-0.35) .. (chem.south east);
\end{tikzpicture}
\end{center}
}

\begin{document}

\packetheading{Lesson 4-5 $\cdot$ Quick $\cdot$ Teacher Edition}{Solving Rational Equations + Chemistry DOK-3 (Q\#2/Q\#6/Q\#7 Prep)}

\sectionbanner{OBJECTIVES}
{\small
\renewcommand{\arraystretch}{1.25}
\noindent\begin{tabular}{|p{1.8in}|p{4.8in}|}
\hline
\lessonrow{Math Objective}{Solve rational equations in one variable by clearing denominators, check for extraneous solutions, and apply rational equations to real-world rate and mixture problems.}
\lessonrow{Language Objective}{Justify why a candidate solution is extraneous, using the domain restrictions of the original rational equation.}
\lessonrow{Essential Understanding}{Multiplying both sides of a rational equation by an expression containing variables can introduce extraneous solutions --- values that satisfy the cleared equation but violate the original domain.}
\end{tabular}
}

\sectionbanner{TOPIC GOALS / ESSENTIAL QUESTION / MATERIALS}
{\small
\renewcommand{\arraystretch}{1.25}
\noindent\begin{tabular}{|p{1.8in}|p{4.8in}|}
\hline
\lessonrow{Topic Goals}{Fluency solving rational equations; identifying extraneous solutions; applying to work-rate and mixture problems (Topic 4 LEHS Q\#2, Q\#6, Q\#7 prep).}
\lessonrow{Essential Question}{When you solve a rational equation by clearing denominators, why must you check every candidate solution against the original equation?}
\lessonrow{Materials}{Student packet, pencil, calculator. DOK-3 Practice \#33 (chemistry mixture performance task) rehearses Topic 4 LEHS Q\#6 (solve rational equation) and Q\#7 (work-rate reasoning).}
\end{tabular}
}

\begin{callout}{\IconPin\ PRE-FLIGHT --- SINGLE-PERIOD COMPRESSED LESSON + DOK-3 CAPSTONE}{calloutpink}
This is a single-period quick treatment of Lesson 4-5. TWO examples are modeled in Launch (deviation from standard Klimsara one-example rule --- solving AND extraneous checking are both assessment-critical). Release promptly at minute 18.\\
Today's capstone is Practice \#33 (chemistry alcohol mixture). Students must (a) set up \(f(x) = (50\cdot0.02 + 0.06x)/(50+x) = 0.05\), (b) clear the denominator and solve --- \(50\cdot0.02 + 0.06x = 0.05(50+x) \rightarrow 1 + 0.06x = 2.5 + 0.05x \rightarrow 0.01x = 1.5 \rightarrow x = 150\) gallons, (c) explain calculator TABLE verification. Same algebraic move as Topic 4 LEHS Q\#6; same reasoning as Q\#7.\\
Pace: Try-It 1 (4 min) \(\rightarrow\) Try-It 3 (4 min) \(\rightarrow\) \#10 (4 min) \(\rightarrow\) \#14 (4 min) \(\rightarrow\) \#25 (6 min) \(\rightarrow\) \#33 (10 min). \(\sim\)32 min total Explore.\\
45-min Wed F variant: cut \#14 + \#25. Never cut \#33.
\end{callout}

\phasetag{bridge from 4-4 to solving rational equations}
\frameworkphaseheader{Do Now}{1}{5}{%
\phaseitems{
  \item Hand out at door.
  \item Circulate 3 min.
  \item Cold-call 1 student on the conflicting candidate solutions.
}%
}{%
\phaseitems{
  \item Read and decide whether either candidate is valid.
  \item Identify which (if any) is extraneous.
}%
}{%
\phaseitems{
  \item What does it mean for a candidate solution to be `extraneous'?
  \item Why do we check candidates in the ORIGINAL equation, not the cleared one?
}%
}{Solo teacher: circulate silently. No hints.}

\begin{callout}{\IconMic\ BRIDGE SCRIPT (read aloud after Do Now)}{calloutblue}
"Topic 4 so far: inverse variation, reciprocal functions, multiplying and adding rational expressions. Today: SOLVING rational EQUATIONS. Same algebra, but with an equation to solve and a new trap --- extraneous solutions. The Model \& Discuss shows two students solving a rational equation and getting different candidate solutions. Which (if any) is valid?"
\end{callout}

\bankitem{Do Now (4-5-savvas-model-discuss-lesson-4-5-launch)}{%
\textbf{CRITIQUE \& EXPLAIN}

Nicky and Tavon used different methods to solve the equation \(\frac{1}{2}x + \frac{2}{5} = \frac{9}{10}\).

\begin{center}
\renewcommand{\arraystretch}{1.25}
\begin{tabular}{|p{2.55in}|p{2.55in}|}
\hline
\textbf{Nicky} & \textbf{Tavon} \\ \hline
\(\frac{1}{2}x + \frac{2}{5} = \frac{9}{10}\) & \(\frac{1}{2}x + \frac{2}{5} = \frac{9}{10}\) \\ \hline
\(\frac{1}{2}x = \frac{9}{10} - \frac{2}{5}\) & \(10\left(\frac{1}{2}x + \frac{2}{5}\right) = 10\left(\frac{9}{10}\right)\) \\ \hline
\(\frac{1}{2}x = \frac{5}{10}\) & \(5x + 4 = 9\) \\ \hline
\(x = 1\) & \(5x = 5\) \\ \hline
The solution is 1. & \(x = 1\) \\ \hline
 & The solution is 1. \\ \hline
\end{tabular}
\end{center}

\begin{enumerate}[label=\textbf{\Alph*.}, leftmargin=1.6em,itemsep=3pt,topsep=3pt]
\item Explain the different strategies that Nicky and Tavon used and the advantages or disadvantages of each.
\item Did Nicky use a correct method to solve the equation? Did Tavon?
\item Use Structure Why might Tavon have chosen to multiply both sides of the equation by 10? Could he have used another number? Explain.
\end{enumerate}
}{}

\teacheranswerbox{\IconCheck\ ANSWER KEY}{%
(verify)
}

\phasetag{teacher models Example 1 + Example 3 [COMPRESSED --- 2 examples]}
\frameworkphaseheader{Launch}{1--2}{13}{%
\phaseitems{
  \item Project Example 1 (solve rational equation). Think aloud: find LCD, multiply both sides, solve polynomial, check in original.
  \item "Every candidate must be checked in the ORIGINAL equation --- clearing denominators can introduce values that don't belong."
  \item Transition to Example 3 without stopping --- announce: "Now watch what happens when a candidate IS extraneous."
  \item Project Example 3. Think aloud: state domain restrictions FIRST, then solve; test each candidate against the restrictions.
  \item Flag the Rules card --- students copy into margin.
  \item 30-sec pause: "Turn to partner: why does clearing denominators sometimes produce extraneous solutions? What's the mechanism?"
  \item Do NOT solve Try-Its or Practice items --- release at minute 18.
}%
}{%
\phaseitems{
  \item Watch Example 1 silently. Note the four-step procedure.
  \item Watch Example 3. Mark the domain-restriction identification step.
  \item During pause: explain mechanism to partner using Rule 1(d).
  \item Copy Rules card into margin.
}%
}{%
\phaseitems{
  \item What is the LCD in Example 1?
  \item After you clear denominators, what kind of equation do you get?
  \item In Example 3, what values are forbidden by the original denominators?
  \item Why does multiplying by an expression containing variables sometimes introduce extraneous roots?
}%
}{Project both examples back-to-back. Release promptly at minute 18.}

\begin{callout}{\IconMic\ TEACHER SCRIPT}{calloutblue}
\begin{enumerate}[leftmargin=1.5em,itemsep=2pt,topsep=2pt]
\item "Watch me solve a rational equation. Step 1: find the LCD."
\item "Step 2: Multiply BOTH SIDES by the LCD. This clears all fractions."
\item "Step 3: Solve the resulting polynomial equation --- you'll get candidate solutions."
\item "Step 4: CHECK each candidate in the ORIGINAL equation. Any candidate that makes a denominator zero is extraneous --- throw it out."
\item "Example 3: watch what happens when a candidate IS extraneous. I state the domain restrictions FIRST --- before I even start solving --- so I know what's forbidden."
\item "Today's DOK-3 task is a chemistry alcohol mixture. Same procedure --- set up rational equation, clear, solve, check. Same algebraic move as Topic 4 LEHS Q\#6."
\item Release to Explore.
\end{enumerate}
\end{callout}

\phasetag{TPS + DOK-3 driver}
\frameworkphaseheader{Explore}{2--3}{32}{%
\phaseitems{
  \item 5 laps across 32 min.
  \item Lap 1 (4 min): Try-It 1 --- solve. Press pairs on the LCD identification.
  \item Lap 2 (4 min): Try-It 3 --- extraneous. Watch for students who skip the check.
  \item Lap 3 (4 min): \#10 --- solve + check. Confirm the check step.
  \item Lap 4 (4 min): \#14 --- extraneous trap. Press "which candidate violates the original domain?" on every pair.
  \item Lap 5 (6 min): \#25 --- work-rate. Use the \(1/a + 1/b = 1/t\) setup.
  \item Lap 6 (10 min): \IconStar\ \#33 Chemistry. Students set up the mixture equation, clear, solve for \(x\), and explain TABLE verification.
  \item Do NOT give \#33 Part A numeric answer. Pairs present argument at Share.
}%
}{%
\phaseitems{
  \item 6 items in order: Try-It 1 \(\rightarrow\) Try-It 3 \(\rightarrow\) \#10 \(\rightarrow\) \#14 \(\rightarrow\) \#25 \(\rightarrow\) \#33.
  \item For \#33: BEFORE computing, write out \(f(x) = (50\cdot0.02 + 0.06x)/(50+x)\) and set equal to 0.05. Ask: what would \(x\) mean in gallons?
  \item Justify TABLE verification with sentence frame.
}%
}{%
\phaseitems{
  \item Try-It 1: what's the LCD of the fractions?
  \item Try-It 3: which candidate makes the original denominator zero?
  \item \#10: after you solve, WHY do you still need to check?
  \item \#14: state the domain restrictions BEFORE you clear denominators.
  \item \#25: how does \(1/a + 1/b = 1/t\) capture the work-rate idea?
  \item \#33 Part A: set up the mixture equation. What does \(x\) represent in gallons --- does your answer make sense?
  \item \#33 Part B: how do you use \(y_1\) and \(y_2\) in the TABLE to verify?
}%
}{Solo teacher: 6 laps. On \#33, press the mixture-equation setup AND the gallons reasonableness check --- those are the assessment-critical moves.}

\bankitem{Try It 1 (4-5-savvas-try-it-1-lesson-4-5)}{%
What is the solution to each equation?

\begin{enumerate}[label=\textbf{\alph*.}, leftmargin=1.6em,itemsep=3pt,topsep=3pt]
\item \(\frac{2}{x+5} = 4\)
\item \(\frac{1}{x-7} = 2\)
\end{enumerate}
}{}

\teacheranswerbox{\IconCheck\ ANSWER / ANTICIPATED TALK}{%
Answer key (from source): \(x = -4.5\)
}

\bankitem{Try It 3 (4-5-savvas-try-it-3-lesson-4-5)}{%
What is the solution to the equation \(\frac{1}{x+2} + \frac{1}{x-2} = \frac{4}{(x+2)(x-2)}\)?
}{}

\teacheranswerbox{\IconCheck\ ANSWER / ANTICIPATED TALK}{%
Answer key (from source): no solution
}

\bankitem{Practice \#10 (4-5-savvas-q10)}{%
Error Analysis Describe and correct the error Miranda made in solving the rational equation.
\[
\begin{aligned}
\frac{1}{x-2} + \frac{x-2}{x+2} &= \frac{x-4}{x-2} \\
(x+2)(x-2)\left(\frac{1}{x-2} + \frac{x-2}{x+2}\right) &= \left(\frac{x-4}{x-2}\right)(x+2)(x-2) \\
(x+2)(1) + (x-2)(x-2) &= (x+4)(x+2) \\
x + 2 + x^2 - 4x + 4 &= x^2 + 6x + 8 \\
-3x + 6 &= 6x + 8 \\
-2 &= 9x; \text{ or } x = -\frac{2}{9}\ \text{X}
\end{aligned}
\]
}{}

\teacheranswerbox{\IconCheck\ ANSWER / ANTICIPATED TALK}{%
Answer key (from source): Miranda incorrectly expanded the right side as \((x+4)(x+2)\) instead of \((x-4)(x+2)\). The correct right side is \(x^2 - 2x - 8\). Solving \(-3x + 6 = -2x - 8\) yields \(x = 14\).
}

\bankitem{Practice \#14 (4-5-savvas-q14)}{%
Make Sense and Persevere Solve the rational equation shown. Explain what is unique about the solution.
\[
\frac{x^2 - 7x - 18}{x+2} = x - 9
\]
}{}

\teacheranswerbox{\IconCheck\ ANSWER / ANTICIPATED TALK}{%
Answer key (from source): The solution is all real numbers except \(x = -2\). Factoring the numerator shows that the equation is an identity \(((x-9)(x+2))/(x+2) = x-9\), which is true for all \(x\) in the domain.
}

\bankitem{Practice \#25 (4-5-savvas-q25)}{%
Make Sense and Persevere Kenji can finish a puzzle in 2 hours working alone. Oscar can finish the same puzzle in 3 hours working alone. How long would it take Oscar and Kenji to finish the puzzle if they worked on it together?
}{}

\teacheranswerbox{\IconCheck\ ANSWER / ANTICIPATED TALK}{%
Answer key (from source): 1.2 hours
}

\begin{callout}{\IconSpeak\ CHAIN 2 CALLBACK --- Chemist Mixture (DE-FACTO PLANT, first rendered item)}{calloutpeach}
\textit{When to say:} Before students start Practice \#33. L35\_P3's q30 (Venetta) was not rendered in that packet, so Chemist becomes the first student-visible Chain 2 item. Plant the archetype here.

\smallskip
``Chain 2 is about \textbf{operating a given model}. Today the chemist hands us a rational equation --- the one that relates alcohol percentage to gallons added --- and we solve for a specific amount. Whereas Chain 1 (Storage Box, Wavelength) asked us to \textbf{build} the model from a constraint, here the equation is \textbf{given}. We operate it. Same move will appear in Half-Life (Topic 5) and Compound Interest (Topic 6) --- different function families each time.''

\smallskip
\textit{Tip:} the extraneous-solution check at the end is shared with Chain 5 (extraneous chain) --- side-note it aloud.
\end{callout}

\bankitem{Practice \#33 \IconStar\ DOK-3 CHEMISTRY MIXTURE (4-5-savvas-q33)}{%
\textbf{Performance Task}

A chemist needs alcohol solution in the correct concentration for her experiment. She adds a 6\% alcohol solution to 50 gallons of solution that is 2\% alcohol. The function that represents the percent of alcohol in the resulting solution is
\[
f(x) = \frac{50(0.02)+x(0.06)}{50+x},
\]
where \(x\) is the amount of 6\% solution added.

\chemistfigure
\placeholdertext{[IMAGE: Chemist pouring liquid into a container; labels ``x gallons of 6\% alcohol'' and ``50 gallons of 2\% alcohol.'']}

\textbf{Part A} How much 6\% solution should be added to create a solution that is 5\% alcohol?

\textbf{Part B} Use Appropriate Tools Explain the steps you could take to use your graphing calculator to verify the correctness of your answer to part (A).
}{}

\begin{callout}{\IconStar\ DOK-3 CHEMISTRY MIXTURE --- Practice \#33 (Topic 4 LEHS Q\#6 + Q\#7 prep)}{calloutpink}
q33 is the chemistry alcohol-mixture performance task: 50 gal of 2\% alcohol solution, add \(x\) gal of 6\% alcohol, resulting concentration \(f(x) = (50\cdot0.02 + 0.06x)/(50+x)\). Part A: solve \(f(x) = 0.05\) for \(x\). Part B: explain how to verify with a graphing calculator table.\\
"The move Part A wants is: set up the rational equation, clear the denominator, solve the resulting linear equation for \(x\). Students should recognize this as a rational equation in disguise --- same moves as Ex 1."\\
For Part B: enter \(y_1 =\) numerator, \(y_2 = 0.05(50+x)\), use TABLE to find where they intersect.\\
If stuck: ask `what would the answer \(x\) mean in gallons?' before they compute, to anchor the reasonableness check.\\
Answer: 150 gallons.\\
Sentence frame required for ELL.
\end{callout}

\teacheranswerbox{\IconCheck\ ANSWER / ANTICIPATED TALK}{%
Answer key (from source): Part A 150 gallons. Part B Using the graphing function of your calculator, enter \(2.5+0.05x\) for \(y_1\), and \(1+0.06x\) for \(y_2\). Then use the TABLE function to find the \(y\)-values of the functions that are equivalent when \(x=150\). \textbullet\ IMAGE PLACEHOLDER [illustration]: Chemist pouring liquid into a container; labels ``x gallons of 6\% alcohol'' and ``50 gallons of 2\% alcohol.''
}

\phasetag{academic conversation + Topic 4 LEHS Q\#2/Q\#6/Q\#7 bridge}
\frameworkphaseheader{Share / Summary}{2}{5}{%
\phaseitems{
  \item Cold-call 1 pair to present their chemistry solution using sentence frame.
  \item Class signals thumbs. Write \(x = 150\) gallons on board.
  \item Topic 4 bridge: "This is the exact move on Topic 4 LEHS Q\#6, the domain reasoning is Q\#2, and the rate reasoning is Q\#7. You have now rehearsed all three."
}%
}{%
\phaseitems{
  \item Presenting pair reads their solution and cites the TABLE verification.
  \item Rest of class signals and writes takeaway in margin.
}%
}{%
\phaseitems{
  \item Summarize: how did you get from \(f(x) = 0.05\) to \(x = 150\)?
  \item Why did 150 gallons make sense as an answer?
}%
}{Cold-call a pair that struggled on \#33 --- hold them accountable to the setup step.}

\phasetag{summary recap (not graded)}
\frameworkphaseheader{Exit Ticket}{1--2}{3}{%
\phaseitems{
  \item Collect at door.
  \item Scan stem 1 (original / extraneous). Plan re-teach if many students can't articulate WHY the check is required.
}%
}{%
\phaseitems{
  \item 3 sentence stems, honest.
}%
}{%
\phaseitems{
  \item Did students connect the `check in original' step to avoiding extraneous solutions?
  \item Did students correctly describe the mixture-equation setup?
}%
}{Collect. No grade.}

\begin{callout}{\IconClipboard\ EXIT TICKET --- Student Form}{calloutyellow}
After clearing denominators in a rational equation, I must check each candidate solution in the \_\_\_\_\_ equation to catch \_\_\_\_\_ solutions.\\
In Practice \#33 (chemistry), the equation came from setting \_\_\_\_\_ equal to the target concentration 0.05.\\
One biggest thing I learned today: \_\_\_\_\_\_\_\_\_\_\_\_\_\_\_\_\_\_\_\_\_\_\_\_\_\_\_\_
\end{callout}

\clearpage

\begin{callout}{\IconClock\ PERIOD A / PERIOD F --- 65-MIN CLASS NOTES}{calloutpink}
Both sections should have 65 min for this lesson. Use the extra 10 min on Explore \#33 (Chemistry) --- richer reasonableness and TABLE-verification talk pays off on Topic 4 LEHS Q\#6/Q\#7.\\
45-min Wed F variant: cut \#14 + \#25 from Explore. Never cut \#33.\\
If finished early: hand out Practice \#8 (work-rate validity reasoning) as fast-finisher; rehearses Q\#7 specifically.
\end{callout}

\begin{callout}{\IconPuzzle\ MODIFICATIONS / SUPPORT FOR IEP STUDENTS}{calloutiep}
\textbullet\ Provide the Solving Rational Equations Rules card as a sticker/cut-out.\\
\textbullet\ For \#33, provide the setup: "\(f(x) = (50\cdot0.02 + 0.06x)/(50+x) = 0.05\). Multiply both sides by \((50+x)\)." Students solve from there.\\
\textbullet\ Accept verbal TABLE-verification explanation in lieu of written steps.
\end{callout}

\begin{callout}{\IconGlobe\ SUPPORTS FOR ELL STUDENTS}{calloutell}
\textbullet\ Sentence frames on student page (don't skip).\\
\textbullet\ Word cards: "extraneous" = an answer that looks right algebraically but doesn't actually work in the original equation; "LCD" = least common denominator; "domain restriction" = a value that \(x\) is not allowed to be.\\
\textbullet\ "Concentration" = amount of substance per unit of mixture; "mixture" = two solutions combined.\\
\textbullet\ Pair ELL students with English-dominant partners for TPS.
\end{callout}

\end{document}
